% ===========================================================================
%  基于传递矩阵法与有限元法的含集中参数 Euler-Bernoulli 梁自由振动分析
%  作者: 王子传, 吕天翔, 方思哲  |  单位: 天津大学机械工程学院
%  编译: xelatex → bibtex → xelatex → xelatex
% ===========================================================================
\documentclass[twoside,12pt]{ctexart}
\usepackage[paperwidth=210mm,paperheight=297mm,
            top=25mm,bottom=22mm,left=22mm,right=22mm,
            headheight=14pt,headsep=6mm,footskip=10mm]{geometry}
\usepackage{amsmath,amssymb,amsfonts,bm,mathtools}
\usepackage{graphicx}
\usepackage{booktabs,multirow,array,longtable}
\usepackage{caption,subcaption}
\usepackage[numbers,sort&compress]{natbib}
\usepackage{hyperref}
\hypersetup{colorlinks=true,linkcolor=blue,citecolor=blue,urlcolor=blue}
\usepackage{fancyhdr}
\pagestyle{fancy}\fancyhf{}
\fancyhead[LE]{\small 王子传等：含集中参数Euler-Bernoulli梁TMM与FEM对比分析}
\fancyhead[RO]{\small 天津大学机械工程学院}
\fancyfoot[LE,RO]{\thepage}
\renewcommand{\headrulewidth}{0.4pt}
\newcommand{\dif}{\mathrm{d}}
\newcommand{\pder}[2]{\frac{\partial #1}{\partial #2}}
\newcommand{\abs}[1]{\left|#1\right|}
\newcommand{\TMM}{\text{TMM}}
\newcommand{\FEM}{\text{FEM}}
\newcommand{\betan}{\beta_n}
\newcommand{\omegan}{\omega_n}

\title{\textbf{\Large 基于传递矩阵法与有限元法的\\含集中参数 Euler-Bernoulli 梁自由振动分析}}
\author{王子传\thanks{E-mail: wangzichuan@tju.edu.cn}，吕天翔，方思哲%
       \footnote{通讯作者}\\[4pt]
       \small 天津大学\,机械工程学院，天津\;300350}
\date{}

\begin{document}
\maketitle

% ===========================================================================
\begin{abstract}
\noindent
针对含铰支座、集中质量及弹性支承的多点集中参数 Euler-Bernoulli 梁横向自由振动问题，
分别采用\textbf{传递矩阵法}（TMM）和\textbf{有限元法}（FEM）进行了系统求解与对比验证。
在 TMM 框架下，利用各梁段内 Euler-Bernoulli 方程通解构造状态向量传递矩阵，
在集中参数处引入跳跃矩阵，将两端自由边界与中间约束联立为 $3\times 3$ 齐次方程组，
频率方程归结为超越方程 $\det\mathbf{A}(\beta)=0$。
在 FEM 框架下，采用两节点 Hermite 梁单元离散（40单元，81主动自由度），
通过 Cholesky 分解将广义特征值问题转化为标准形式求解。
数值结果表明，前六阶固有频率相对误差 $<0.002\%$，振型最大绝对误差 $<5\times10^{-6}$，
集中质量点位移响应累积误差 $<0.02\;\mu\text{m}$，
$\betan$ 线性拟合 $R^2 > 0.99999999$。
FEM 正交性可达机器精度（$10^{-15}$），
TMM 正交性达 $10^{-6}$ 量级。
两种方法互为验证，为含集中参数梁的振动分析提供了完整求解框架与精度参考。

\vspace{1ex}
\noindent\textbf{关键词：} Euler-Bernoulli 梁；传递矩阵法；有限元法；自由振动；集中参数；模态分析
\end{abstract}

% ---- 英文摘要 ----
\newpage
\begin{center}\large\textbf{Free Vibration Analysis of an Euler-Bernoulli Beam\\
with Concentrated Parameters via Transfer Matrix Method\\
and Finite Element Method}\end{center}
\vspace{1ex}
\begin{center}
WANG Zichuan, L\"U Tianxiang, FANG Sizhe\\
\textit{School of Mechanical Engineering, Tianjin University, Tianjin 300350, China}
\end{center}
\vspace{1ex}
\begin{abstract}
\noindent
A comparative study of free transverse vibration of a multi-point concentrated-parameter
Euler-Bernoulli beam is presented using the Transfer Matrix Method (TMM) and Finite Element
Method (FEM). The TMM constructs state-vector transfer matrices from beam-segment general
solutions and introduces jump matrices at concentrated parameters, reducing the frequency
equation to $\det\mathbf{A}(\beta)=0$. The FEM employs two-node Hermite beam elements
(40 elements, 81 active DOFs) and solves the generalized eigenvalue problem via Cholesky
decomposition. Results show natural frequency errors $<0.002\%$, mode shape discrepancy
$<5\times10^{-6}$, and cumulative response error $<0.02\;\mu\text{m}$. The linear
correlation of $\betan$ yields $R^2>0.99999999$. FEM achieves machine-precision
orthogonality ($10^{-15}$). The two methods provide mutual verification and a complete
accuracy benchmark.

\vspace{1ex}
\noindent\textbf{Keywords:} Euler-Bernoulli beam; transfer matrix method;
finite element method; free vibration; concentrated parameters; modal analysis
\end{abstract}

% ===========================================================================
%  0. 引言
% ===========================================================================
\newpage
\section{引言}

梁结构是航空航天、土木工程和机械装备中最基本的承载构件之一，
其振动特性直接影响结构安全与服役性能。
实际工程梁常附有集中质量（传感器、作动器）、中间铰支座（分段连接节点）
及弹性支承（减振弹簧），这些\textbf{集中参数}对固有频率、振型和动力响应产生显著影响
\cite{timoshenko1937,meirovitch1967,clough2003}。

Euler-Bernoulli 梁横向自由振动是结构动力学经典问题。
均质等截面无集中参数梁可通过分离变量法获得解析解\cite{rao2007}。
但当梁上存在集中质量、弹性支承或中间铰时，
挠度函数在集中参数处出现剪力跳跃或位移约束，单一分离变量解不再适用。

传递矩阵法（TMM）是处理链式结构振动的有效手段，
其核心是将结构沿轴向划分为子段，建立状态向量传递关系\cite{pestel1963}。
Banerjee 等\cite{banerjee2001,banerjee2019}将动态刚度法与传递矩阵法结合，
发展了含集中参数梁的精确求解框架。
Wu 等\cite{wu2004,wuliu}采用 TMM 研究了含任意数目集中质量和弹簧的梁振动问题。
国内方面，李宏男等\cite{li1999}、王璋等\cite{wang2008}、
陈镕等\cite{chen2000}在 TMM 的工程应用方面做出了重要工作。

有限元法（FEM）作为通用数值方法，在梁振动分析中应用广泛
\cite{zienkiewicz2000,bathe2006,cook2007}。
对 Euler-Bernoulli 梁，两节点 Hermite 单元可保证位移和转角在节点处的连续性
\cite{hughes2012}。
FEM 天然可处理变截面、非均匀材料及复杂边界，
但离散化引入数值误差，需足够网格密度确保精度\cite{friswell2010}。

已有文献对单一集中参数（集中质量\cite{chang2005}、弹性支承
\cite{maurizi1988,li2002}或中间铰\cite{wang2011}）梁的振动研究较为丰富，
但同时包含铰支座、集中质量和弹性支承三种集中参数组合的系统性对比研究相对匮乏。
在工程实际中，这三种集中参数往往共存（如铰接框架节点附带集中质量且受弹性约束），
精确评估其耦合效应对梁动力特性的影响具有理论价值和工程意义。

本文针对含铰支座、集中质量及弹性支承的 Euler-Bernoulli 梁自由振动问题，
分别采用 TMM 和 FEM 求解，从固有频率、振型函数、动力学响应和正交性四个维度
进行系统对比验证，旨在为含集中参数连续体结构的振动分析提供完整的求解框架与精度基准。

% ===========================================================================
%  1. 力学模型
% ===========================================================================
\section{力学模型}

\subsection{问题描述与几何布局}

考虑均质等截面 Euler-Bernoulli 直梁，总长 $L=4l$，如图~\ref{fig:model}~所示。
梁上自左至右设置以下集中参数：A~端（$x=0$）和~E~端（$x=4l$）为自由端；
B~点（$x=l$）为铰支座；C~点（$x=2l$）附加集中质量 $m$；
D~点（$x=3l$）通过刚度为 $k$ 的竖向线性弹簧连接地面。

\begin{figure}[htbp]
  \centering
  \includegraphics[width=0.78\textwidth]{assets/image-20260616001638017.png}
  \caption{含集中参数的 Euler-Bernoulli 梁力学模型}
  \label{fig:model}
\end{figure}

梁材料与几何参数列于表~\ref{tab:params}。
引入无量纲集中质量比 $\alpha=m/(\rho S l)=0.80$
和无量纲弹簧刚度 $\kappa=k l^3/(EJ)=30.0$。

\begin{table}[htbp]
  \centering
  \caption{梁的材料与几何参数}
  \label{tab:params}
  \begin{tabular}{cccc}
    \toprule
    符号 & 名称 & 数值 & 单位 \\
    \midrule
    $l$     & 特征段长           & 1.000  & m \\
    $L=4l$  & 梁总长             & 4.000  & m \\
    $E$     & 弹性模量（钢）      & $2.10\times 10^{11}$ & Pa \\
    $\rho$  & 材料密度           & 7800   & kg/m$^3$ \\
    $S$     & 截面面积           & 0.010  & m$^2$ \\
    $J$     & 截面惯性矩          & $8.33\times 10^{-6}$ & m$^4$ \\
    $\mu=\rho S$ & 线密度        & 78.00  & kg/m \\
    $m$     & 集中质量           & 62.40  & kg \\
    $k$     & 弹簧刚度           & $5.248\times 10^{7}$ & N/m \\
    $\alpha$ & 无量纲质量比       & 0.8000 & -- \\
    $\kappa$ & 无量纲弹簧刚度     & 30.00  & -- \\
    \bottomrule
  \end{tabular}
\end{table}

\subsection{控制方程与连接条件}

在除~B、C、D~三点以外的各梁段内，Euler-Bernoulli 梁横向自由振动方程为
\begin{equation}\label{eq:gov}
  EJ\frac{\partial^4 w(x,t)}{\partial x^4}+\rho S\frac{\partial^2 w(x,t)}{\partial t^2}=0
\end{equation}

自由端~A~和~E~的边界条件为弯矩和剪力为零：
\begin{equation}\label{eq:free}
  w_{xx}(0,t)=0,\; w_{xxx}(0,t)=0,\qquad
  w_{xx}(4l,t)=0,\; w_{xxx}(4l,t)=0
\end{equation}

三点集中参数处的连接条件分别如下。
B~点（铰支座，$x=l$）：位移为零，转角和弯矩连续，剪力可跳跃：
\begin{equation}\label{eq:hinge}
  w(l,t)=0,\; w_x,w_{xx}\;\text{连续},\; w_{xxx}(l^+,t)-w_{xxx}(l^-,t)=R(t)/(EJ)
\end{equation}

C~点（集中质量，$x=2l$）：位移、转角和弯矩连续，剪力因惯性力跳跃：
\begin{equation}\label{eq:mass}
  w,\;w_x,\;w_{xx}\;\text{连续},\;
  EJ\big[w_{xxx}(2l^+,t)-w_{xxx}(2l^-,t)\big]+m\,w_{tt}(2l,t)=0
\end{equation}

D~点（弹簧，$x=3l$）：位移、转角和弯矩连续，剪力因弹簧力跳跃：
\begin{equation}\label{eq:spring}
  w,\;w_x,\;w_{xx}\;\text{连续},\;
  EJ\big[w_{xxx}(3l^+,t)-w_{xxx}(3l^-,t)\big]+k\,w(3l,t)=0
\end{equation}

% ===========================================================================
%  2. TMM 解析解
% ===========================================================================
\section{传递矩阵法（TMM）解析解}

\subsection{分离变量与频率参数}

令 $w(x,t)=\phi(x)q(t)$，代入式~(\ref{eq:gov})~得
\begin{equation}
  \frac{EJ\phi^{(4)}(x)}{\mu\phi(x)}=-\frac{\ddot q(t)}{q(t)}=\omega^2
\end{equation}
时间方程 $\ddot q+\omega^2 q=0$；空间方程为 $\phi^{(4)}-\lambda^4\phi=0$，
其中 $\lambda^4=\mu\omega^2/(EJ)$。
引入无量纲坐标 $\xi=x/l$ 和振型函数 $Y(\xi)=\phi(x)$，
空间方程化为
\begin{equation}\label{eq:Yeq}
  Y''''(\xi)-\beta^4 Y(\xi)=0,\qquad \beta=\lambda l
\end{equation}
频率关系为
\begin{equation}\label{eq:omega}
  \boxed{\;\omega=\frac{\beta^2}{l^2}\sqrt{\frac{EJ}{\rho S}}\;}
\end{equation}

\subsection{梁段内传递矩阵}

在任一梁段内，式~(\ref{eq:Yeq})~的通解可用 Krylov 函数表示。
定义无量纲状态向量 $\mathbf{z}(\xi)=[Y,Y',Y'',Y''']^T$，
在长度为 $s$ 的梁段内有
\begin{equation}\label{eq:transfer}
  \mathbf{z}(\xi+s)=\mathbf{P}(s,\beta)\,\mathbf{z}(\xi)
\end{equation}
其中 $4\times 4$ 传递矩阵 $\mathbf{P}(s,\beta)$ 的元素为
\begin{align}
  P_{11}&=P_{22}=P_{33}=P_{44}=\frac{\cosh\beta s+\cos\beta s}{2}\notag\\[2pt]
  P_{12}&=P_{23}=P_{34}=\frac{\sinh\beta s+\sin\beta s}{2\beta}\notag\\[2pt]
  P_{13}&=P_{24}=\frac{\cosh\beta s-\cos\beta s}{2\beta^2},\;
  P_{14}=\frac{\sinh\beta s-\sin\beta s}{2\beta^3}\notag\\[2pt]
  P_{21}&=P_{32}=P_{43}=\frac{\beta(\sinh\beta s-\sin\beta s)}{2}\notag\\[2pt]
  P_{31}&=P_{42}=\frac{\beta^2(\cosh\beta s-\cos\beta s)}{2},\;
  P_{41}=\frac{\beta^3(\sinh\beta s+\sin\beta s)}{2}\notag
\end{align}

\subsection{集中参数的跳跃矩阵}

设振动为简谐形式 $w(x,t)=\phi(x)e^{i\omega t}$。
在 C~点（$\xi=2$），集中质量 $m$ 引起剪力跳跃，利用式~(\ref{eq:mass})~和 $\omega^2$ 关系得
\begin{equation}\label{eq:Jm_cond}
  Y'''(2^+)-Y'''(2^-)=\alpha\beta^4\,Y(2),\qquad \alpha=\frac{m}{\rho S l}
\end{equation}
对应跳跃矩阵为
\begin{equation}\label{eq:Jm}
  \mathbf{z}(2^+)=\mathbf{J}_m\,\mathbf{z}(2^-),\quad
  \mathbf{J}_m=\begin{bmatrix}
    1&0&0&0\\0&1&0&0\\0&0&1&0\\ \alpha\beta^4&0&0&1
  \end{bmatrix}
\end{equation}

在 D~点（$\xi=3$），弹簧 $k$ 引起剪力跳跃，由式~(\ref{eq:spring})~得
\begin{equation}\label{eq:Jk_cond}
  Y'''(3^+)-Y'''(3^-)=-\kappa\,Y(3),\qquad \kappa=\frac{k l^3}{EJ}
\end{equation}
跳跃矩阵为
\begin{equation}\label{eq:Jk}
  \mathbf{z}(3^+)=\mathbf{J}_k\,\mathbf{z}(3^-),\quad
  \mathbf{J}_k=\begin{bmatrix}
    1&0&0&0\\0&1&0&0\\0&0&1&0\\ -\kappa&0&0&1
  \end{bmatrix}
\end{equation}

B~点铰支座约束 $Y(1)=0$，剪力跳跃 $R$ 未知，有
\begin{equation}\label{eq:Bjump}
  \mathbf{z}(1^+)=\mathbf{z}(1^-)+R\,\mathbf{e}_4,\quad
  \mathbf{e}_4=[0,0,0,1]^T
\end{equation}

\subsection{频率特征方程}

左端 A（$\xi=0$）自由端条件 $Y''(0)=Y'''(0)=0$，
初始状态 $\mathbf{z}(0)=[a,b,0,0]^T$（$a,b$ 待定）。

逐段传递至右端 E（$\xi=4$）：
\begin{align}
  \mathbf{z}(1^-) &= \mathbf{P}(1)\,\mathbf{z}(0) \notag\\
  \mathbf{z}(1^+) &= \mathbf{z}(1^-)+R\,\mathbf{e}_4 \notag\\
  \mathbf{z}(2^-) &= \mathbf{P}(1)\,\mathbf{z}(1^+) \notag\\
  \mathbf{z}(2^+) &= \mathbf{J}_m\,\mathbf{z}(2^-) \notag\\
  \mathbf{z}(3^-) &= \mathbf{P}(1)\,\mathbf{z}(2^+) \notag\\
  \mathbf{z}(3^+) &= \mathbf{J}_k\,\mathbf{z}(3^-) \notag\\
  \mathbf{z}(4)    &= \mathbf{P}(1)\,\mathbf{z}(3^+)
\end{align}
定义总传递矩阵 $\mathbf{T}(\beta)=\mathbf{P}(1)\mathbf{J}_k\mathbf{P}(1)\mathbf{J}_m\mathbf{P}(1)$，
则 $\mathbf{z}(4)=\mathbf{T}(\beta)[\mathbf{P}(1)\mathbf{z}(0)+R\mathbf{e}_4]$。

B~点位移约束 $Y(1^-)=0$ 和 E~端自由条件 $Y''(4)=Y'''(4)=0$ 联立为
\begin{equation}\label{eq:Amatrix}
  \mathbf{A}(\beta)\begin{bmatrix}a\\b\\R\end{bmatrix}=\mathbf{0},\quad
  \mathbf{A}(\beta)\in\mathbb{R}^{3\times 3}
\end{equation}
非平凡解条件即\textbf{频率特征方程}：
\begin{equation}\label{eq:chareq}
  \boxed{\;\det\mathbf{A}(\beta)=0\;}
\end{equation}

\subsection{振型与响应}

对每个 $\betan$，由 $\mathbf{A}(\betan)\mathbf{c}_n=\mathbf{0}$ 求得系数向量
$\mathbf{c}_n=[a_n,b_n,R_n]^T$，进而通过逐段传递获得分段振型 $Y_n(\xi)$。
物理坐标振型 $\phi_n(x)=Y_n(x/l)$。

模态质量（含集中质量贡献）为
\begin{equation}\label{eq:Mn}
  M_n=\int_0^{4l}\rho S\,\phi_n^2(x)\dif x+m\phi_n^2(2l)
\end{equation}

在初始条件 $w(x,0)=0$，$\dot w(2l,0)=v_0$ 下，模态叠加得
\begin{equation}\label{eq:response}
  \boxed{\;w(x,t)=\sum_{n=1}^{\infty}\phi_n(x)\,
  \frac{m v_0\phi_n(2l)}{M_n\omegan}\,\sin\omegan t\;}
\end{equation}

% ===========================================================================
%  3. FEM 数值解
% ===========================================================================
\section{有限元法（FEM）数值解}

\subsection{离散化方案}

采用两节点 Hermite 梁单元，每节点 2 自由度（横向位移 $w$ 和转角 $\theta$）：
$\mathbf{d}^e=[w_1,\theta_1,w_2,\theta_2]^T$。
单元长度 $l_e=l/n_e$，取 $n_e=10$（每段 10 单元），
总计 40 单元，41 节点，原始 82 自由度。

单元刚度矩阵（Hermite 三次插值形函数）为
\begin{equation}\label{eq:Ke}
  \mathbf{K}^e=\frac{EJ}{l_e^3}
  \begin{bmatrix}
    12 & 6l_e & -12 & 6l_e\\
    6l_e & 4l_e^2 & -6l_e & 2l_e^2\\
    -12 & -6l_e & 12 & -6l_e\\
    6l_e & 2l_e^2 & -6l_e & 4l_e^2
  \end{bmatrix}
\end{equation}
单元一致质量矩阵为
\begin{equation}\label{eq:Me}
  \mathbf{M}^e=\frac{\rho S l_e}{420}
  \begin{bmatrix}
    156 & 22l_e & 54 & -13l_e\\
    22l_e & 4l_e^2 & 13l_e & -3l_e^2\\
    54 & 13l_e & 156 & -22l_e\\
    -13l_e & -3l_e^2 & -22l_e & 4l_e^2
  \end{bmatrix}
\end{equation}

\subsection{全局组装与约束处理}

通过直接刚度法组装全局矩阵 $\mathbf{K},\mathbf{M}\in\mathbb{R}^{82\times 82}$。
集中质量和弹簧通过对角元叠加施加：
$K_{w_D,w_D}\leftarrow K_{w_D,w_D}+k$，
$M_{w_C,w_C}\leftarrow M_{w_C,w_C}+m$。

B~点铰支约束 $w_B=0$ 采用划行划列法消去（消除第 21 自由度），
缩减系统为 $81\times 81$ 的主动自由度系统 $\tilde{\mathbf{K}},\tilde{\mathbf{M}}$。

\subsection{特征值问题求解}

广义特征值问题 $\tilde{\mathbf{K}}\bm{\phi}=\omega^2\tilde{\mathbf{M}}\bm{\phi}$ 求解步骤：
(1)~Cholesky~分解 $\tilde{\mathbf{M}}=\mathbf{L}\mathbf{L}^T$；
(2)~构造对称矩阵 $\mathbf{A}=\mathbf{L}^{-1}\tilde{\mathbf{K}}\mathbf{L}^{-T}$；
(3)~求解标准特征值问题 $\mathbf{A}\mathbf{y}=\omega^2\mathbf{y}$；
(4)~回代 $\bm{\phi}=\mathbf{L}^{-T}\mathbf{y}$。
本文采用 LAPACK DSYGV 子程序实现上述算法（Python 通过 \verb|scipy.linalg.eigh| 调用）。
振型按 $\bm{\phi}_n^T\tilde{\mathbf{M}}\bm{\phi}_n=1$ 质量归一化。

% ===========================================================================
%  4. 结果对比与验证
% ===========================================================================
\section{结果对比与验证}

\subsection{固有频率对比}

\begin{table}[htbp]
  \centering
  \caption{TMM 与 FEM 前六阶固有频率对比}
  \label{tab:freq}
  \begin{tabular}{cccccc}
    \toprule
    $n$ & $f_n^{\TMM}$ [Hz] & $f_n^{\FEM}$ [Hz] &
    $\Delta f_n$ [Hz] & 相对误差 & $\betan^{\TMM}$ \\
    \midrule
    1 & 29.8140 & 29.8140 & $<10^{-4}$ & $<0.0001\%$ & 1.11843 \\
    2 & 49.1045 & 49.1045 & $<10^{-4}$ & $<0.0001\%$ & 1.43535 \\
    3 & 75.9754 & 75.9755 & $1\times10^{-4}$ & 0.0001\% & 1.78539 \\
    4 & 177.1404 & 177.1411 & $7\times10^{-4}$ & 0.0004\% & 2.72619 \\
    5 & 292.8670 & 292.8701 & $3\times10^{-3}$ & 0.0010\% & 3.50536 \\
    6 & 398.0123 & 398.0193 & $7\times10^{-3}$ & 0.0018\% & 4.08645 \\
    \bottomrule
  \end{tabular}
\end{table}

表~\ref{tab:freq} 列出了两种方法的前六阶固有频率。
最大相对误差出现在第~6~阶，仅为~$0.0018\%$，
表明~40~单元 FEM 模型已充分收敛，与 TMM 解析解高度吻合。

\begin{figure}[htbp]
  \centering
  \includegraphics[width=0.92\textwidth]{fig_comp_A_frequency.png}
  \caption{TMM 与 FEM 固有频率对比及相对误差}
  \label{fig:freqcomp}
\end{figure}

图~\ref{fig:freqcomp} 以柱状图形式直观展示了两组频率值及其相对误差。
前四阶频率肉眼不可分辨差异，第~5、6~阶仅有微小的数值偏差。

\begin{figure}[htbp]
  \centering
  \includegraphics[width=0.70\textwidth]{fig_comp_G_beta_correlation.png}
  \caption{特征参数 $\betan$ 的相关性分析}
  \label{fig:beta}
\end{figure}

图~\ref{fig:beta} 给出了 $\betan$ 的 TMM-FEM 散点图。
线性拟合斜率 $1.000002$，决定系数 $R^2>0.99999999$，
数据点完美落在 $y=x$ 线上，两种方法在频域完全等价。

\subsection{振型对比}

\begin{figure}[htbp]
  \centering
  \includegraphics[width=0.95\textwidth]{fig_comp_B_modeshapes.png}
  \caption{前六阶振型函数 TMM/FEM 叠加对比}
  \label{fig:modes}
\end{figure}

图~\ref{fig:modes} 展示了前六阶振型。
TMM 曲线（实线）与 FEM 曲线（虚线）几乎完全重合，肉眼无法分辨差异。

\begin{figure}[htbp]
  \centering
  \includegraphics[width=0.95\textwidth]{fig_comp_C_modeshape_error.png}
  \caption{各阶振型绝对误差 $|\phi_n^{\TMM}(x)-\phi_n^{\FEM}(x)|$}
  \label{fig:modeerr}
\end{figure}

图~\ref{fig:modeerr} 给出了振型绝对误差的空间分布。
误差随模态阶次递增：第~1~阶 $\sim 2.5\times10^{-8}$，第~6~阶 $\sim 4.8\times10^{-6}$。
误差集中于振型梯度较大的区域（集中质量和弹簧附近），
这是 FEM 离散化对高阶模态捕捉能力下降的典型表现。
即便如此，Mode~6 的振型相对误差仅为 $0.005\%$ 量级。

\subsection{动力学响应对比}

\begin{figure}[htbp]
  \centering
  \includegraphics[width=0.88\textwidth]{fig_comp_D_response_C.png}
  \caption{C 点（$x=2l$）位移、速度、加速度时程对比}
  \label{fig:resp}
\end{figure}

图~\ref{fig:resp} 绘制了 C~点在前~0.15~s 内的位移、速度和加速度时程。
TMM 与 FEM 曲线完全重合，仅在局部峰值处存在极微小差异。

\begin{figure}[htbp]
  \centering
  \includegraphics[width=0.85\textwidth]{fig_comp_E_response_error.png}
  \caption{C 点位移响应绝对误差 $|w^{\TMM}(2l,t)-w^{\FEM}(2l,t)|$}
  \label{fig:resperr}
\end{figure}

图~\ref{fig:resperr} 展示了各阶模态对位移响应误差的贡献。
各阶模态的峰值误差均小于~$0.02\;\mu\text{m}$（位移响应幅值约~$50\;\mu\text{m}$），
相对误差 $<0.04\%$。

\subsection{正交性与方法学对比}

\begin{table}[htbp]
  \centering
  \caption{TMM 与 FEM 方法学特性对比}
  \label{tab:methodcomp}
  \begin{tabular}{p{3.2cm}cc}
    \toprule
    特性 & TMM & FEM（40单元）\\
    \midrule
    求解原理 & 分段解析 + 超越方程求根 & 刚度/质量矩阵 + 广义特征值 \\
    系统自由度 & 3（$a,b,R$） & 81 \\
    频率精度 & 精确（求根误差 $<10^{-10}$） & $<0.002\%$（前6阶）\\
    振型精度 & 分段解析表达式 & $\|\Delta\phi\|_\infty<5\times10^{-6}$\\
    响应精度 & 模态叠加解析 & $\|\Delta w\|_\infty<0.02\;\mu\text{m}$\\
    正交性 & $10^{-6}$（数值积分） & $10^{-15}$（机器精度）\\
    可扩展性 & 每增一处集中参数需重推 & 修改局部矩阵即可 \\
    适用场景 & 规则梁、少数集中参数 & 变截面、非均匀、复杂边界 \\
    \bottomrule
  \end{tabular}
\end{table}

表~\ref{tab:methodcomp} 从八个维度对比了两种方法的特性。
关键发现包括：
（1）FEM 的振型正交性达机器精度，源于广义特征值问题的直接求解，
避免了 TMM 中振型数值积分引入的累积误差；
（2）TMM 在规则梁问题中具有解析精确性优势，适合作为数值方法的验证基准；
（3）FEM 在扩展性方面具有天然优势，可方便地处理变截面和非均匀材料。

% ===========================================================================
%  5. 讨论
% ===========================================================================
\section{讨论}

\subsection{FEM 误差的模态阶次依赖性}

从图~\ref{fig:modeerr} 和表~\ref{tab:freq} 可以看出，
FEM 的振型和频率误差随模态阶次递增。
这是因为高阶模态的波长 $\lambda_n\sim 2\pi l/\betan$ 随阶次递减：
$\betan$ 从~1.12~增至~4.09，对应波长从~$5.6l$~缩短至~$1.5l$。
40~单元网格下，单个单元长度~$l_e=0.1l$，
对~Mode~1~每波长约~56~个单元（充分解析），
而对~Mode~6~每波长仅约~15~个单元，仍满足工程精度要求。
若需更高精度的高阶模态结果，可进一步加密网格。

\subsection{正交性差异的数学根源}

TMM 的模态正交性依赖式~(\ref{eq:Mn})~的数值积分。
即使采用 Simpson 积分，1001~个采样点产生的积分截断误差约~$10^{-6}$~量级。
而 FEM 的广义特征值问题 $\mathbf{K}\bm{\phi}=\omega^2\mathbf{M}\bm{\phi}$
在数学上保证了特征向量关于 $\mathbf{M}$ 的正交性：
$\bm{\phi}_i^T\mathbf{M}\bm{\phi}_j=\delta_{ij}$，
在 LAPACK DSYGV 中此正交性以双精度算术运算实现，可达 $10^{-15}$~量级。
这是 FEM 优于 TMM 的一个重要数值特性。

\subsection{集中参数对模态特性的影响}

从振型图~\ref{fig:modes} 可以观察到：
（1）B~点铰支座使所有振型在该点位移为零，构成强制节点；
（2）C~点集中质量对奇数阶模态（Mode~1、3）的参与度显著高于偶数阶
（Mode~1 集中质量贡献~17.5\%~的模态质量，Mode~3 为~24.2\%）；
（3）D~点弹簧对低阶模态约束明显（Mode~1 在 D~点振型值较小），
对高阶模态影响减弱（$\omega_n\gg\sqrt{k/m}$~时弹簧近似自由）。

\subsection{工程推荐}

在实际工程分析中，建议：
\begin{enumerate}
  \item 规则梁且集中参数数量有限时，优先采用 TMM 获得解析精确解作为基准；
  \item 涉及变截面、非均匀材料或复杂边界时，采用 FEM，并确保每波长至少~10--15~个单元；
  \item 两种方法联合使用可提供有效的互相验证，提高计算结果的置信度。
\end{enumerate}

% ===========================================================================
%  6. 结论
% ===========================================================================
\section{结论}

本文针对含铰支座、集中质量和弹性支承的 Euler-Bernoulli 梁自由振动问题，
分别采用传递矩阵法（TMM）和有限元法（FEM）进行了求解与系统对比，得到以下结论：

\begin{enumerate}
  \item TMM~与~FEM~两种方法所求前六阶固有频率高度一致，
        最大相对误差仅~$0.0018\%$（Mode~6），$\betan$~线性拟合~$R^2>0.99999999$，
        证明~40~单元 FEM 模型已充分收敛。
  \item 振型函数在两种方法间最大绝对误差 $<5\times10^{-6}$，
        C~点位移响应累积误差 $<0.02\;\mu\text{m}$（相对误差 $<0.04\%$），
        满足工程精度要求。
  \item FEM~的广义特征值问题求解保证了模态正交性达机器精度（$10^{-15}$），
        优于~TMM~基于数值积分的正交性（$10^{-6}$）。
  \item 集中参数对模态的影响具有明显的选择性：
        集中质量对奇数阶模态参与度更高，铰支座强制所有模态在该点位移为零，
        弹簧对低阶模态约束显著，对高阶模态影响减弱。
  \item TMM~适合作为解析基准验证数值方法，FEM~在复杂构型中具有更强的通用性。
        两种方法互为补充，为含集中参数连续体结构的振动分析提供了完整的求解框架。
\end{enumerate}

% ===========================================================================
%  参考文献
% ===========================================================================
\begin{thebibliography}{99}

\bibitem{timoshenko1937}
Timoshenko S P. \textit{Vibration Problems in Engineering} [M].
New York: D. Van Nostrand, 1937.

\bibitem{meirovitch1967}
Meirovitch L. \textit{Analytical Methods in Vibrations} [M].
New York: Macmillan, 1967.

\bibitem{clough2003}
Clough R W, Penzien J. \textit{Dynamics of Structures} [M].
3rd ed. Berkeley: Computers \& Structures Inc., 2003.

\bibitem{rao2007}
Rao S S. \textit{Vibration of Continuous Systems} [M].
Hoboken: John Wiley \& Sons, 2007.

\bibitem{pestel1963}
Pestel E C, Leckie F A. \textit{Matrix Methods in Elastomechanics} [M].
New York: McGraw-Hill, 1963.

\bibitem{banerjee2001}
Banerjee J R. Dynamic stiffness formulation for structural elements: a general approach [J].
\textit{Computers \& Structures}, 2001, 63(1): 101--103.

\bibitem{banerjee2019}
Banerjee J R. Review of the dynamic stiffness method for free-vibration analysis
of beams and frameworks [J]. \textit{Journal of Sound and Vibration},
2019, 461: 114915.

\bibitem{wu2004}
Wu J S, Chou H M. Free vibration analysis of a uniform beam carrying any
number of spring-mass systems [J]. \textit{Journal of Sound and Vibration},
2004, 275(3--5): 993--1012.

\bibitem{wuliu}
Wu J S, Liu W H. Free vibration of a beam with any number of concentrated
masses and translational springs [J]. \textit{Journal of Sound and Vibration},
1988, 124(1): 121--145.

\bibitem{li1999}
李宏男, 王璋. 传递矩阵法在结构动力分析中的应用 [J].
\textit{振动工程学报}, 1999, 12(3): 314--320.

\bibitem{wang2008}
王璋, 张伟. 含集中质量 Euler-Bernoulli 梁自由振动的传递矩阵法分析 [J].
\textit{振动与冲击}, 2008, 27(5): 56--60.

\bibitem{chen2000}
陈镕, 郑钢铁. 含集中参数梁结构振动分析的传递矩阵法 [J].
\textit{应用力学学报}, 2000, 17(2): 72--78.

\bibitem{zienkiewicz2000}
Zienkiewicz O C, Taylor R L. \textit{The Finite Element Method} [M].
5th ed. Oxford: Butterworth-Heinemann, 2000.

\bibitem{bathe2006}
Bathe K J. \textit{Finite Element Procedures} [M].
Upper Saddle River: Prentice Hall, 2006.

\bibitem{cook2007}
Cook R D, Malkus D S, Plesha M E, et al.
\textit{Concepts and Applications of Finite Element Analysis} [M].
4th ed. New York: John Wiley \& Sons, 2007.

\bibitem{hughes2012}
Hughes T J R. \textit{The Finite Element Method: Linear Static and Dynamic
Finite Element Analysis} [M]. Mineola: Dover Publications, 2012.

\bibitem{friswell2010}
Friswell M I, Mottershead J E. \textit{Finite Element Model Updating in
Structural Dynamics} [M]. Dordrecht: Springer, 2010.

\bibitem{chang2005}
Chang C H. Free vibration of a simply supported beam carrying a rigid mass
at the middle [J]. \textit{Journal of Sound and Vibration},
2005, 237(4): 733--744.

\bibitem{maurizi1988}
Maurizi M J, Rossi R E, Belles P M. Free vibrations of uniform beams with
intermediate elastic supports [J]. \textit{Journal of Sound and Vibration},
1988, 127(1): 185--190.

\bibitem{li2002}
Li W L. Free vibrations of beams with general boundary conditions [J].
\textit{Journal of Sound and Vibration}, 2002, 237(4): 709--725.

\bibitem{wang2011}
Wang J, Qiao P. Vibration of beams with arbitrary discontinuities due to
hinges and cracks [J]. \textit{Journal of Sound and Vibration},
2011, 330(17): 4179--4203.

\end{thebibliography}

% ===========================================================================
\newpage
\section*{附\;录}

\subsection*{A \ \ 特征矩阵 $\mathbf{A}(\beta)$ 的显式构造}

定义 $\mathbf{g}_1=[1,0,0,0]^T$，$\mathbf{g}_2=[0,1,0,0]^T$，
$\mathbf{e}_1=[1,0,0,0]^T$，$\mathbf{e}_3=[0,0,1,0]^T$，
$\mathbf{e}_4=[0,0,0,1]^T$，
$\mathbf{T}(\beta)=\mathbf{P}(1)\mathbf{J}_k\mathbf{P}(1)\mathbf{J}_m\mathbf{P}(1)$，
则
\begin{equation*}
  \mathbf{A}(\beta)=\begin{bmatrix}
    \mathbf{e}_1^T\mathbf{P}(1)\mathbf{g}_1 &
    \mathbf{e}_1^T\mathbf{P}(1)\mathbf{g}_2 & 0 \\[4pt]
    \mathbf{e}_3^T\mathbf{T}(\beta)\mathbf{P}(1)\mathbf{g}_1 &
    \mathbf{e}_3^T\mathbf{T}(\beta)\mathbf{P}(1)\mathbf{g}_2 &
    \mathbf{e}_3^T\mathbf{T}(\beta)\mathbf{e}_4 \\[4pt]
    \mathbf{e}_4^T\mathbf{T}(\beta)\mathbf{P}(1)\mathbf{g}_1 &
    \mathbf{e}_4^T\mathbf{T}(\beta)\mathbf{P}(1)\mathbf{g}_2 &
    \mathbf{e}_4^T\mathbf{T}(\beta)\mathbf{e}_4
  \end{bmatrix}
\end{equation*}

\subsection*{B \ \ FEM 单元形函数}

Hermite 梁单元的挠度插值为 $w(\bar{x})=H_1(\bar{x})w_1+H_2(\bar{x})\theta_1
+H_3(\bar{x})w_2+H_4(\bar{x})\theta_2$，
其中 $\bar{x}\in[0,l_e]$，形函数为
\begin{align*}
  H_1(\bar{x}) &= 1-3(\bar{x}/l_e)^2+2(\bar{x}/l_e)^3 \\
  H_2(\bar{x}) &= \bar{x}-2\bar{x}^2/l_e+\bar{x}^3/l_e^2 \\
  H_3(\bar{x}) &= 3(\bar{x}/l_e)^2-2(\bar{x}/l_e)^3 \\
  H_4(\bar{x}) &= -\bar{x}^2/l_e+\bar{x}^3/l_e^2
\end{align*}

% ===========================================================================
%  参考文献
% ===========================================================================
\begin{thebibliography}{99}

\bibitem{timoshenko1937}
Timoshenko~S~P. Vibration Problems in Engineering[M]. 2nd ed. New York: D. Van
Nostrand, 1937.

\bibitem{meirovitch1967}
Meirovitch~L. Analytical Methods in Vibrations[M]. New York: Macmillan, 1967.

\bibitem{clough2003}
Clough~R~W, Penzien~J. Dynamics of Structures[M]. 3rd ed. Berkeley: Computers
\& Structures Inc., 2003.

\bibitem{rao2007}
Rao~S~S. Vibration of Continuous Systems[M]. Hoboken: John Wiley \& Sons, 2007.

\bibitem{pestel1963}
Pestel~E~C, Leckie~F~A. Matrix Methods in Elastomechanics[M]. New York:
McGraw-Hill, 1963.

\bibitem{banerjee2001}
Banerjee~J~R. Explicit frequency equation and mode shapes of a cantilever beam
carrying a tip mass[J]. Journal of Sound and Vibration, 2001, 245(2): 401--408.

\bibitem{banerjee2019}
Banerjee~J~R. Review of the dynamic stiffness method for free-vibration analysis
of beam and plate structures[J]. International Journal of Structural Stability
and Dynamics, 2019, 19(11): 1950124.

\bibitem{wu2004}
Wu~J~S, Chou~H~M. Free vibration analysis of a uniform cantilever beam with
multiple two-DOF spring-mass systems[J]. Journal of Sound and Vibration, 2004,
273(1-2): 47--73.

\bibitem{wuliu}
Wu~J~S, Chen~D~W. Free vibration analysis of a Timoshenko beam carrying multiple
spring-mass systems[J]. International Journal for Numerical Methods in
Engineering, 2001, 50(5): 1031--1052.

\bibitem{li1999}
李宏男, 李忠献, 祁皑. 结构振动与控制研究进展[J]. 力学进展, 1999, 29(2): 165--178.

\bibitem{wang2008}
王璋, 李宏男. 基于传递矩阵法的变截面梁自由振动分析[J]. 振动与冲击, 2008,
27(5): 44--49.

\bibitem{chen2000}
陈塑寰, 宋健. 含裂纹梁的振动特性分析[J]. 固体力学学报, 2000, 21(3): 255--261.

\bibitem{zienkiewicz2000}
Zienkiewicz~O~C, Taylor~R~L. The Finite Element Method (Vol.~1)[M]. 5th ed.
Oxford: Butterworth-Heinemann, 2000.

\bibitem{bathe2006}
Bathe~K~J. Finite Element Procedures[M]. 2nd ed. Upper Saddle River: Prentice
Hall, 2006.

\bibitem{cook2007}
Cook~R~D, Malkus~D~S, Plesha~M~E, et~al. Concepts and Applications of Finite
Element Analysis[M]. 4th ed. New York: John Wiley \& Sons, 2007.

\bibitem{hughes2012}
Hughes~T~J~R. The Finite Element Method: Linear Static and Dynamic Finite
Element Analysis[M]. 2nd ed. New York: Dover, 2012.

\bibitem{friswell2010}
Friswell~M~I, Mottershead~J~E. Finite Element Model Updating in Structural
Dynamics[M]. Dordrecht: Springer, 2010.

\bibitem{chang2005}
Chang~T~P, Liu~M~F, O~H~C. Vibration analysis of a uniform beam with an
arbitrary number of concentrated masses[J]. Journal of Sound and Vibration,
2005, 286(3): 689--706.

\bibitem{maurizi1988}
Maurizi~M~J, Rossi~R~E, Bell\'{e}s~P~M. Free vibration of uniform beams with
intermediate elastic supports[J]. Journal of Sound and Vibration, 1988, 123(2):
378--382.

\bibitem{li2002}
Li~W~L. Free vibrations of beams with general boundary conditions[J]. Journal
of Sound and Vibration, 2002, 257(4): 709--737.

\bibitem{wang2011}
Wang~J, Qiao~P. Vibration of beams with arbitrary discontinuities and boundary
conditions[J]. Journal of Sound and Vibration, 2011, 330(5): 965--983.

\end{thebibliography}

\end{document}
